\chapter{Construire et lancer le programme}
\label{chap-utilisation}

\section{Prérequis}

Un compilateur C11, CMake~3.25 ou plus récent, \code{pkg-config} et les fichiers de
développement de GTK~3. GStreamer est facultatif (effets sonores).

\begin{table}[ht]
    \centering
    \small
    \begin{tabularx}{\textwidth}{@{}l>{\raggedright\arraybackslash}X@{}}
        \toprule
        \textbf{Système} & \textbf{Installation des dépendances}                               \\
        \midrule
        Debian, Ubuntu   & \code{sudo apt install build-essential cmake pkg-config libgtk-3-dev
                               libgstreamer1.0-dev} \\
        Fedora           & \code{sudo dnf install gcc cmake pkgconf-pkg-config gtk3-devel
                               gstreamer1-devel}       \\
        macOS            & \code{brew install cmake pkg-config gtk+3 gstreamer}                \\
        \bottomrule
    \end{tabularx}
\end{table}

\section{Commandes}

\begin{lstlisting}[language=bash, morekeywords={make, cmake, ctest, git}]
git clone https://github.com/DmytroPalahin/battleship.git
cd battleship

make test        # configurer, compiler, lancer les tests
make run         # lancer le jeu
make play        # jouer dans le terminal
make simulate    # mesurer les trois niveaux (10 000 parties)
make asan        # tests sous AddressSanitizer et UBSan
make report      # compiler ce rapport
make help        # liste complète
\end{lstlisting}

Le \fichier{Makefile} n'est qu'un aide-mémoire : il appelle les préréglages CMake. La forme
explicite est \code{cmake -{}-preset release}, \code{cmake -{}-build -{}-preset release},
\code{ctest -{}-preset release}.

\section{Options de ligne de commande}

\begin{table}[ht]
    \centering
    \small
    \begin{tabularx}{\textwidth}{@{}l>{\raggedright\arraybackslash}X@{}}
        \toprule
        \multicolumn{2}{@{}l}{\textbf{\code{battleship}} --- application graphique}               \\
        \midrule
        \code{-l}, \code{-{}-level=NIVEAU} & Démarrer directement à un niveau : \code{easy},
        \code{medium} ou \code{hard}.                                                             \\
        \code{-s}, \code{-{}-seed=N}       & Graine de la première partie, pour la rejouer.       \\
        \code{-m}, \code{-{}-mute}         & Désactiver les effets sonores.                       \\
        \code{-V}, \code{-{}-version}      & Afficher la version.                                 \\
        \code{-{}-autoplay}                & Test : l'ordinateur joue les deux camps, puis le
        programme se ferme.                                                                       \\
        \code{-{}-max-shots=N}             & Test : avec \code{-{}-autoplay}, s'arrêter après
        $N$ tirs.                                                                                 \\
        \code{-{}-screenshot=FICHIER}      & Test : enregistrer la fenêtre au format PNG avant de
        quitter.                                                                                  \\
        \addlinespace
        \multicolumn{2}{@{}l}{\textbf{\code{battleship-cli}} --- interface terminal}              \\
        \midrule
        \code{play}                        & Jouer contre l'ordinateur (\code{-{}-level},
        \code{-{}-seed}).                                                                         \\
        \code{simulate}                    & Mesurer les niveaux (\code{-{}-games},
        \code{-{}-seed}, \code{-{}-level}, \code{-{}-csv},
        \code{-{}-no-duel}).                                                                      \\
        \bottomrule
    \end{tabularx}
\end{table}

\section{Commandes du jeu}

À la souris : cliquer sur une case de la flotte pirate. Au clavier : les flèches déplacent
le missile sur la grille adverse, \touche{Espace} tire, \touche{F11} bascule en plein écran.
